\section{Background and Threat Model}\label{sec:background}

\subsection{ZK applications}\label{sec:bg-apps}

A ZK application lets a user prove a statement about private data and have a
contract act on it. Examples include withdrawing from a privacy pool, minting
to the holder of a credential, voting once per identity, settling a rollup
batch, and accepting a zkVM execution result. Such an application has three
parts.

\paragraph{Circuit}
The circuit, or zkVM guest program, defines a relation $R(\vec{u},\vec{w})$
over public inputs $\vec{u}$ and private witnesses $\vec{w}$. It is written
in a circuit language (Circom, Noir, gnark, halo2) or as an ordinary program
for a zkVM (RISC Zero, SP1). A proof $\pi$ convinces a verifier that the
prover knows some $\vec{w}$ with $R(\vec{u},\vec{w})=1$, and it reveals
nothing else.

\paragraph{Verifier}
The verifier is usually generated from the circuit's verification key and
deployed as a contract (Groth16~\cite{groth16}, PLONK~\cite{plonk}), or reached
through a router or an aggregation layer. It takes $(\vec{u},\pi)$ and returns
accept or reject. It knows nothing about what $\vec{u}$ means.

\paragraph{Consumer}
The consumer is the contract that calls the verifier and, on acceptance,
commits effects. It assembles $\vec{u}$ from calldata and storage, so the
consumer decides which of its values are covered by the proof. Every other
value it uses comes from the transaction sender.

The public-input vector is the only interface between the three parts. Each
circuit and each consumer defines its own meaning for the entries. There is no
type system, no shared schema, and usually no test that exercises a
malicious caller with an honest proof.

\subsection{Threat model}\label{sec:bg-threat}

\paragraph{Attacker}
The attacker is any account that can send transactions. It controls all
calldata, its own address, and the value it sends, and it can observe pending
and past transactions, including other users' valid proofs. It can resubmit an
observed proof with different arguments and can try to be included first.

\paragraph{Assumptions}
The attacker cannot break the cryptographic assumptions of the proof system or
the hash functions, and it cannot forge proofs for false statements of a sound
verifier. Administrators and deployers are trusted. Findings that require a
malicious owner are out of scope.

\paragraph{Goal}
The attacker obtains an effect that the proved statement does not authorize:
funds or rights sent to itself, a larger amount, a repeated effect, an effect
in a context the statement was not made for, or a change to another user's
state.

\paragraph{Out of scope}
Soundness bugs inside a proof system (transcript, polynomial commitment, or
recursion constraints), circuit under-constraints that do not change what the
consumer is told, prover-side bugs, trusted-setup issues, and denial of
service against one's own transaction. Prior work and existing tools target
the first two classes~\cite{sok-snark-bugs,picus,circomspect}. We study the
layer above them.
